\documentclass{../base/canon}

%% canon_variables.tex — Valores por defecto de metadatos institucionales
%%
%% Incluir ANTES del \begin{document} para sobreescribir los defaults de la clase.
%% El template puede sobreescribir cualquier valor después de este \input.
%%
%% Uso en template:
%%   %% canon_variables.tex — Valores por defecto de metadatos institucionales
%%
%% Incluir ANTES del \begin{document} para sobreescribir los defaults de la clase.
%% El template puede sobreescribir cualquier valor después de este \input.
%%
%% Uso en template:
%%   %% canon_variables.tex — Valores por defecto de metadatos institucionales
%%
%% Incluir ANTES del \begin{document} para sobreescribir los defaults de la clase.
%% El template puede sobreescribir cualquier valor después de este \input.
%%
%% Uso en template:
%%   \input{../base/canon_variables}
%%   \SetDocTitle{Mi Informe Específico}   % sobreescribe el default

\SetDocLogo{../assets/logo.pdf}
\SetDocUnit{Tecnología \& Sistemas}
\SetContactUnit{Mesa de soporte}
\SetContactMembers{%
  Equipo CoreX & Soporte & soporte@empresa.com%
}
\SetContactNote{Para soporte técnico o consultas sobre este documento.}

%%   \SetDocTitle{Mi Informe Específico}   % sobreescribe el default

\SetDocLogo{../assets/logo.pdf}
\SetDocUnit{Tecnología \& Sistemas}
\SetContactUnit{Mesa de soporte}
\SetContactMembers{%
  Equipo CoreX & Soporte & soporte@empresa.com%
}
\SetContactNote{Para soporte técnico o consultas sobre este documento.}

%%   \SetDocTitle{Mi Informe Específico}   % sobreescribe el default

\SetDocLogo{../assets/logo.pdf}
\SetDocUnit{Tecnología \& Sistemas}
\SetContactUnit{Mesa de soporte}
\SetContactMembers{%
  Equipo CoreX & Soporte & soporte@empresa.com%
}
\SetContactNote{Para soporte técnico o consultas sobre este documento.}

\SetDocType{PLAN TÉCNICO}
\SetDocTitle{Plan Técnico — Nombre del Proyecto}
\SetDocCode{PLN-TEC-000}
\SetDocArea{Tecnología \& Sistemas}
\SetDocAuthor{Arquitecto / Responsable}
\SetDocDate{\today}

\begin{document}

\section{Justificación}

\begin{ParrafoEjecutivo}
Describir el problema o necesidad que motiva este plan técnico y el beneficio esperado.
\end{ParrafoEjecutivo}

\section{Alcance}

\begin{itemize}
  \item Lo que \textbf{incluye} este plan.
  \item Lo que \textbf{no incluye} (fuera de alcance).
\end{itemize}

\section{Diseño técnico}

\begin{ParrafoMetodologico}
Describir la arquitectura de la solución: componentes involucrados, flujo de datos,
integraciones y decisiones de diseño relevantes.
\end{ParrafoMetodologico}

\subsection{Diagrama o esquema}

\FiguraPlaceholder[0.85\linewidth]{Diagrama de arquitectura — pendiente}

\section{Fases de implementación}

\begin{table}[H]
  \centering
  \caption{Cronograma por fases}
  \begin{adjustbox}{max width=\linewidth}
  \begin{tabular}{l l l l r}
    \toprule
    \textbf{Fase} & \textbf{Descripción} & \textbf{Inicio} & \textbf{Fin} & \textbf{Días} \\
    \midrule
    Fase 1 & Análisis y diseño         & 2026-05-01 & 2026-05-07  & 7  \\
    Fase 2 & Desarrollo                & 2026-05-08 & 2026-05-21  & 14 \\
    Fase 3 & Pruebas y validación      & 2026-05-22 & 2026-05-28  & 7  \\
    Fase 4 & Despliegue y monitoreo    & 2026-05-29 & 2026-06-04  & 7  \\
    \bottomrule
  \end{tabular}
  \end{adjustbox}
\end{table}

\section{Riesgos técnicos}

\WarningBox[Riesgos identificados]{
  Describir los riesgos técnicos de mayor impacto y sus planes de mitigación.
}

\section{Criterios de aceptación}

\begin{enumerate}
  \item Criterio técnico de éxito 1 — con métrica medible.
  \item Criterio técnico de éxito 2 — con métrica medible.
  \item Sin regresiones en funcionalidad existente.
\end{enumerate}

\SignatureBlock{Nombre Apellido}{Responsable Técnico}

\ContactSignature

\end{document}
